\section{课程主线：Frontier Product Research 不是“研究做完再包装”}

这堂课的标题容易被误读成一节普通的强化学习（Reinforcement Learning, RL）导论。Karina Nguyen 真正讨论的是更具体的系统问题：当模型进入写作、软件、协作、教育和长期任务后，\textbf{产品希望模型怎样行动}、\textbf{研究怎样把这种行动写成环境与奖励}、\textbf{界面怎样让用户暴露失败}，必须在同一个闭环里共同设计。若研究先优化一个脱离产品的 benchmark，再把模型交给产品团队包装，很多真实失败直到上线才会出现。

\lecturefigure{slide-01-golden-age-curiosity.jpg}{“好奇、知识创造与制作”的黄金时代：讲座先给出价值方向}{00:00:49--00:01:41}

\teachervoice{00:00:49--00:01:38，Karina 明确说这不是单向 lecture，而希望课堂协作讨论；她的起点是每个人都能借助 AI 建造有意义的东西。这个价值判断解释了为什么后文同时讨论 creativity、agency、verification 与 alignment。}

\begin{importantbox}{定义：Frontier Product Research}
本讲中的 \term{frontier product research} 不是“追踪最新模型并做产品”。它把尚未稳定的模型能力当作研究材料，把真实用户任务构造成训练环境与评测，再用产品交互快速收集反例。其最小闭环是
\[
\begin{aligned}
\text{product belief}&\rightarrow \text{task/env}\rightarrow \text{reward/eval}\\
&\rightarrow \text{model behavior}\rightarrow \text{user evidence}\rightarrow \text{next hypothesis}.
\end{aligned}
\]
其中 product belief 是希望产品最终呈现的行为原则；它必须被翻译成可观察任务，而不能停留在口号。
\end{importantbox}

\begin{knowledgebox}{本讲的证据边界}
前半段大量使用产品 demo。这些画面能证明某个交互在当时可运行，能展示讲者如何理解产品方向，却不能单独证明模型在所有用户、分布和风险等级上可靠。后文的 eval、拒答 taxonomy、XSTest、RL environment 与 reward hacking，正是在补齐 demo 无法回答的问题。
\end{knowledgebox}

\subsection{Canvas：界面改变不是装饰，而是任务结构改变}

第一个例子把教育请求拆成解释、代码生成、执行和可视化。若系统只提供一条聊天消息，用户很难同时阅读概念、修改代码和观察图形；Canvas 把“对话”升级为共享工作区，使模型输出成为可编辑对象。读图时不要只看 Gaussian 曲线，而要看左侧解释、中央代码和右侧执行结果怎样构成一个可验证循环。

\lecturefigure{slide-02-gaussian-canvas.jpg}{Gaussian 分布解释、Python 代码与运行结果被放进同一工作区}{00:02:05--00:02:57}

对随机变量 $X\sim\mathcal{N}(\mu,\sigma^2)$，密度为
\[
p(x)=\frac{1}{\sqrt{2\pi}\sigma}
\exp\!\left(-\frac{(x-\mu)^2}{2\sigma^2}\right).
\]
这里 $\mu$ 控制中心，$\sigma$ 控制尺度。模型给出公式并不等于教学完成；只有当用户能修改参数、运行代码并观察曲线变化时，抽象符号才变成可操作的理解。产品设计因此改变了评测问题：不再只问答案是否正确，还要问代码能否执行、图形是否对应解释、修改是否局部生效。

\teachervoice{00:02:05--00:04:04，讲者用 Gaussian 教学和论文截图解释 Canvas 的动机：ChatGPT 最初是 conversational UI，但代码、长文和细粒度追问暴露了聊天界面的限制。课堂提示：UI 变化来自任务结构，而不是视觉改版。}

论文理解例子进一步展示“验证入口”。用户先提供论文中的 calibration 图，再要求模型解释、生成研究报告和代码。原始材料、模型解释与新产物必须并排可追溯，否则系统容易把顺畅叙述误当成事实。

\lecturefigure{slide-03-paper-calibration-source.jpg}{论文中的 calibration 图：模型解释必须回到原始视觉证据}{00:02:57--00:03:09}

\lecturefigure{slide-04-calibration-canvas-report.jpg}{Canvas 中生成的 calibration 报告：选择局部内容后继续追问}{00:03:09--00:04:12}

\begin{knowledgebox}{首次使用：Calibration}
\term{Calibration（校准）}描述模型置信度与实际正确率是否一致。若所有置信度约为 $q$ 的预测中，真实正确比例也约为 $q$，则该区间近似校准。常见误区是把 calibration 与 accuracy 混为一谈：模型可以准确但过度自信，也可以不够准确却在概率意义上校准。产品需要把“答案是什么”和“应相信多少”分开显示。
\end{knowledgebox}

\lecturefigure{slide-05-claude-usage-bachelor-degrees.jpg}{教育使用报告中的相关性图：产品数据可以产生研究问题}{00:04:12--00:04:45}

这张使用图把 Claude.ai 的教育用途与美国不同本科专业联系起来。读图先看横轴专业占比或使用占比，再看散点偏离趋势线的位置；它能提示不同学科采用模式不同，却不能证明“专业训练导致使用”或“使用导致学习效果”。产品 telemetry 可以产生值得研究的假设，但相关性、选择偏差与用户自报必须单独处理。

\begin{warningbox}{从产品日志到科学结论的三道门}
第一，使用者不是总体的随机样本；第二，点击或生成次数不等于任务成功；第三，群体相关性不能直接解释个体因果。Frontier product research 可以从真实流量快速发现问题，但仍需要明确 sampling、outcome 与 counterfactual。
\end{warningbox}

\subsection{本章小结}

课程开场已经给出整堂课的尺度：模型能力、工作界面、用户证据与后训练目标是一个系统。Canvas 的价值不是“多一个窗口”，而是把解释、编辑、执行和验证放进同一交互；产品数据的价值不是天然真实，而是把研究问题暴露得更具体。

